%% Copyright 2009 Jeffrey D. Hein
%
% This work may be distributed and/or modified under the
% conditions of the LaTeX Project Public License, either version 1.3
% of this license or (at your option) any later version.
% The latest version of this license is in
%   http://www.latex-project.org/lppl.txt
% and version 1.3 or later is part of all distributions of LaTeX
% version 2005/12/01 or later.
%
% This work has the LPPL maintenance status `maintained'.
% 
% The Current Maintainer of this work is Jeffrey D. Hein.
%
% This work consists of the files 3dplot.sty and 3dplot.tex

%Description
%-----------
%3dplot.tex - an example file demonstrating the use of the 3dplot.sty package.

%Created 2009-11-07 by Jeff Hein.  Last updated: 2009-11-09
%----------------------------------------------------------

%Update Notes
%------------

%2009-11-07: Created file along with 3dplot.sty package


\documentclass{article}
%%%<
\usepackage{verbatim}
%%%>

\usepackage{tikz}   %TikZ is required for this to work.  Make sure this exists before the next line

\usepackage{tikz-3dplot} %requires 3dplot.sty to be in same directory, or in your LaTeX installation

\usepackage[active,tightpage]{preview}  %generates a tightly fitting border around the work
\PreviewEnvironment{tikzpicture}
\setlength\PreviewBorder{2mm}

\begin{document}

%Angle Definitions
%-----------------

%set the plot display orientation
%synatax: \tdplotsetdisplay{\theta_d}{\phi_d}
\tdplotsetmaincoords{70}{110}

%define polar coordinates for some vector
%TODO: look into using 3d spherical coordinate system
\pgfmathsetmacro{\rvec}{.8}
\pgfmathsetmacro{\thetavec}{30}
\pgfmathsetmacro{\phivec}{60}

%start tikz picture, and use the tdplot_main_coords style to implement the display 
%coordinate transformation provided by 3dplot
\begin{tikzpicture}[scale=3,tdplot_main_coords,
cube/.style={very thick,fill=green},
angle/.style={red,thick},
curve/.style={thick}]

\tikzstyle{every node}=[font=\huge]

%set up some coordinates 
%-----------------------
\coordinate (O) at (0,0,0);

%determine a coordinate (P) using (r,\theta,\phi) coordinates.  This command
%also determines (Pxy), (Pxz), and (Pyz): the xy-, xz-, and yz-projections
%of the point (P).
%syntax: \tdplotsetcoord{Coordinate name without parentheses}{r}{\theta}{\phi}
\tdplotsetcoord{P}{\rvec}{\thetavec}{\phivec}

%draw figure contents
%--------------------

%draw the main coordinate system axes
\draw[thick,->] (0,0,0) -- (3,0,0) node[anchor=north east]{$z$};
\draw[dotted,->] (-3,0,0) -- (0,0,0);
\draw[thick,->] (0,0,0) -- (0,3,0) node[anchor=north west]{$x$};
\draw[dotted,->] (0,-3,0) -- (0,0,0);
\draw[thick,->] (0,0,0) -- (0,0,3) node[anchor=south]{$y$};
\draw[dotted,->] (0,0,-3) -- (0,0,0);


\draw[cube, fill opacity=.2] (0,-2.5,-2.5) -- (0,-2.5,2.5) -- (0,2.5,2.5) -- (0,2.5,-2.5) -- cycle;

\draw[thick,->] (0,0,0) -- (0,0,1) node[anchor=west]{$p_2$};
\draw[thick,->] (0,0,0) -- (0,1,0) node[anchor=south]{$p_1$};

\draw[dashed] (-3,2,0) -- (0,2,0) node[anchor=north west]{};
\draw[thick,->] (0,2,0) -- (0.9,2,0) node at (0.5,1.6,0) {$2 \cdot t_1$};

\draw[dashed] (-3,0,2) -- (-0.4,0,2) node[anchor=north west]{};
\draw[thick,->] (-0.4,0,2) -- (0.9,0,2) node at (1.0,0.45,2.1){$2 \cdot t_2$};


%define the rotated coordinate frame based on the angle
\tdplotsetthetaplanecoords{90}
		
%draw the circle in the rotated frame
\tdplotdrawarc[curve,tdplot_rotated_coords]{(O)}{2}{0}{360}{}{}

%Bob
\draw plot [mark=*, mark size=1,color=black] coordinates{(0,0,0)} node[above left] {Bob}; 

%Star 1
\draw plot [mark=*, mark size=1,color=black] coordinates{(0,2,0)} node at(-0.5,2.2,0) {$(2,0,0)$}; 
\draw plot [mark=*, mark size=1,color=black] coordinates{(1,2,0)} node at (1.5,1.7,0) {$(2,0,1)$};

%Star 2
\draw plot [mark=*, mark size=1,color=black] coordinates{(-0.4,0,2)} node at (-0.4,0.7,2.05) {$(0,2,-0.4)$}; 
\draw plot [mark=*, mark size=1,color=black] coordinates{(1,0,2)} node at(1.4,-0.3,2) {$(0,2,1)$}; 



\end{tikzpicture}

\end{document}